\chapter{Duplicity} % 10

\vspace{-1.3cm}
\begin{localsize}{10}
    \begin{quote}
        “What you do speaks so loudly that I cannot hear what you say.”
        \begin{flushright}- Ralph Waldo Emerson \end{flushright}
    \end{quote}
\end{localsize}
\vspace{1cm}


The Old Oak stood eerily grey in the otherwise lush Glade of Representative. Its stiff leafless branches swayed naked in warm late spring breeze. The massive trunk, now hollowed out and riddled with newly excavated apartments, had sustained too much damage to pump nutrients to the canopy. Rot had began to toke hold, and larvae of all kinds feasted blindly on its dying roots.

But to Eaglewing, perched high on a barren branch, rot was simply the price of progress.
He looked down upon the Glade with a profound swelling pride. \textit{``We've done it!"}, he thought to himself. The vote had passed just over a month ago, and his many plans to better the Forest was in full motion. He peeked up to the sky, \textit{``Look father, look have far we've come."}

Below him, the quite glade was broken by the heavy, rhythmic panting of three sturdy roebucks. They were dragging a massive, rectangular slab of limestone across the grass.

``Where do you want this?" the largest roebuck grunted, catching his breath as the trio stopped to rest.

``At the opposite end of the Glade," Eaglewing answered, his voice commanding, channeling authority. ``And raise it on its short side, so it stands tall."

``And, how do you expect us to do that? the roebuck responded irritably. Eaglewing gave him a stern look, and the roebuck quickly added, ``Nevermind, we'll figure it out." The bucks lower their heads, and continued to haul the stone.

Eaglewing watched them go with an amused, sly grin. \textit{``Try to forget me now"}, he thought as his mind drifted to his childhood, and everything he had endured. \\

Eaglewing had not always been known as Eaglewing. He had been born Óttarr, named after his father's great-grandfather. Even as a hatchling, his left wing---a dull, patchy colour of brown---stuck out like a sore thumb. A group of older, cruel fledglings had swiftly dubbed him ``Mudwing".

The nickname stuck, as ill-intentioned names tend to do. By his third winter, Óttarr had had enough. He went to his father, Ongenþeow, then the Twig, and wept over the cruelty of his peers.

His father had not comforted him. Instead, he offered cold, hard truth.

``Nicknames, especially those dispensed with malice, are like quicksand---the more you squirm, the deeper they sink." Ongenþeow had said, without even looking at his son. ``The solution, I reckon, is no different than to any other problem. It is solely an affair of shaping public oppinion. You must never directly oppose your peers; float along like a log in the river and wait for an opportunity. And, when opportunity arise, steer it in the direction of your chosing. When you learn how to control rivers, there is no lanscape you cannot mold, and no opposition you cannot erode."

Eager to gain his father's approval and to change his name, the young crow began to hatch a plan.\\

The first week he simply stopped taking offense when called Mudwing. The second week, he pretended to enjoy it. By the third week, he was loudly preaching to anyone who would listen about the magnificent maroon hues of mud, acting deliberately pretentious and high on himself. He was setting bait, giving the bullies a reason to knock him down a notch. He was moving the needle at a speed below perception.

He enlisted the help of his only friend, Brightbeak, keeping him in the dark about the true nature of the scheme. When the day finally came, Ottarr and Brightbeak were walking along the very fledglings who had tormented him.

``Eagles are the rats of heaven," Ottarr declared, unprovoked.

``No, they are not," snapped the largest of the bullies. ``They are majestic. I've even seen one myself."

``I doubt that," Ottarr scoffed. ``I you had, you would know how incredibly ugly they are, with their dull, brown wings."

``But Mudwing," Brightbeak chimed in innocently, right on cue. ``Your wing is brown."

``No, my wing is a brilliant maroon!" Ottarr cried, feigning deep offense.

``No it isn't. It's the exact same color as the eagle I saw," the bully sneered slyly, falling for the bait. ``You have an eagle's wing."

``How dare you compare me with one of those flying rats!" Ottarr shrieked, acting devastated.

``I think we know what this calls for," another bully laughted. The large fledgling caught on, grinning maliciously. ``Surely we do. Don't we, \textit{Eaglewing}?"

Ottar let out a wail of despair. He tucked his head under his wing and fled. Behind him, the bullies roared with laughter. They had found---or so they thought---his new weakness.

For weeks, they relentlessly called him Eaglewing. The yelled it across the glades to ridicule him; they whispered it when he walked past. Ottarr played hurt every time, scurrying away in shame.

But though the bullies were spreading the name with malice, the rest of the murder didn't understand the nature of the joke. When they heard \textit{Eaglewing}, they heard strength. When the younger females heard it, they heard valor. The name began to circulate through the branches outside of the group of bullies, completely stripped of its cruel context.

By the time the bullies realized that their name had backfired---it was too late. The public had adopted the name, and not even the bullies could change it now. When the bullies later went to confront Ottarr about his deceit, Ottarr simply thank them for their service, and flew away.

Ottarr won, and he would never have to be Mudwing ever again.\\

(Scene 3 - Eaglewing sending for a woodpecker

%\input{testcuneiform.tex}